\section{Scope, source audit and the added mechanisms}\label{sec:scope}

The source of this investigation is the ProveIt directory
\path{Logic/PrimitiveRecursiveArithmetic}. We fixed the snapshot
\begin{center}
\code{d56cd5505fb93135b90b048bf3dcca9e73615162}.
\end{center}
Its index lists 57 research packets, including six in \code{foundations/}.
It distinguishes these delivered packets from the checked ConciseIPRA
library and says that the packets were filed without rerunning their
verification commands~\cite{repository}. That distinction governs this
report: source claims are evidence of what a packet proposes, not automatic
imports of its mathematical correctness. We inspected the foundation
packets and the particular analytic and Diophantine reports relevant below;
we have not audited all 57 packets.

The target theory is first-order intuitionistic logic with equality, the
defining equations for every fixed primitive-recursive definition, successor
injectivity and nonzeroness, and induction for quantifier-free formulas.
Parameters in induction are allowed. This is the theory represented by
\code{PRAAx} in the current ConciseIPRA source. Its existing
\code{pra\_iff\_native} theorem transfers derivations to the small native
scalar presentation~\cite{kernel,bridge}. The theorem is a transport map;
it does not create an arithmetic derivation from an ordinary Lean theorem.

\subsection{What is inherited and what is added}

The earlier packets already explain Booleanization of bounded formulas,
finite observations, explicit approximation moduli, precision propagation,
finite-window stability, simple implicit branches and several Runge-type
integer-point arguments~\cite{oldframework,embeddingpacket,finiteobservation,
evidencepacket,uniformpacket,frameworkdated}. Repeating those descriptions
would not answer the present request well. The following table makes the
additional work explicit.

\begin{longtable}{@{}>{\raggedright\arraybackslash}p{.20\linewidth}>{\raggedright\arraybackslash}p{.34\linewidth}>{\raggedright\arraybackslash}p{.40\linewidth}@{}}
\toprule
Interface & Existing coverage & Addition in this report \\
\midrule\endhead
Finite proofs & Quantifier-free induction and finite contracts; several
semantic Lean lemmas and thin object-theory adapters & A proved Lean
transformer for arbitrary primitive recursion with state, time and
parameters, returning a \code{Prf PRAAx} derivation \\
Classical reasoning & Strong emphasis on constructive presentation & A
precise criterion allowing classical first-order PRA proofs of universal
decidable conclusions, with a written negative-translation proof \\
Analysis & Fixed-function constructions, integration, implicit branches,
explicit moduli & A uniform rational-polynomial ODE constructor,
finite refinement and uniqueness bounds, and polynomial conservation \\
Integral points & Square and higher-power gaps; conditional height descent &
A separated polynomial-factor certificate for curves with mixed terms,
explicit complete enumeration, and a uniform exact example family \\
\bottomrule
\end{longtable}

These are additions to the inspected collection. We make no claim of
historical priority for negative translation, Euler's method, discrete
Gronwall estimates, conservation laws or Runge's method. The contribution is
their explicit organization for this weak first-order setting, the proofs
of the stated concrete versions, and the implemented object-theory rule.

\subsection{Four separate questions}

An embedding should answer four questions in order.
\begin{enumerate}
\item \emph{What is represented?} Specify finite codes, validity conditions,
equivalence of representatives and the exact meaning of approximations.
\item \emph{What is computed?} Give fixed PR definitions and numerical bounds,
including every search bound, precision demand and iteration count.
\item \emph{Why is it correct in the weak theory?} Supply a finite invariant
and show that its defining assertion is quantifier-free, or give its
proved Booleanization.
\item \emph{What has actually been formalized?} Distinguish a mathematical
proof, an ordinary Lean proof and a derivation in the embedded theory.
\end{enumerate}

For example, a record with a field asserting a global height bound does
not produce a height bound. A type-correct theorem target does not prove
that target. Similarly, defining a rational Euler algorithm does not prove
that its approximation errors converge. The work below supplies the latter
mathematical proofs and implements one general part of the fourth question.

\subsection{The common local-to-global pattern}

The construction has a finite state and numerical challenge. A local step
has a proved arithmetic contract. The state stores enough history or
accumulated evidence to express a quantifier-free invariant. Primitive
recursion propagates that invariant to a challenge-dependent stopping time.
Finally, a quantitative estimate converts the finite result into the desired
represented conclusion.

For an ODE, the challenge is accuracy, the state is the finite Euler table,
and the decisive estimate compares two finite meshes. For the integral-point
application, the challenge is an encoded candidate point; separated factors
force a numerical height bound, after which a finite table lists all answers.
Both use the same arithmetic proof discipline, although their mathematical
estimates are different. The framework is useful precisely because those
estimates remain visible instead of being hidden inside a universal
``certificate soundness'' assumption.
